\documentclass[11pt]{article}
\usepackage{mathblatt}
\def\mitzone{1}
\begin{document}
\blattkopf{Funktionsklassen und Eigenschaften}{Gesamt}
\verzeichniszeile{\verz{zone}{Kennst du schon (Nr.~1–8)} \verztrenn \verz{e1}{1 Funktionswert und Punkt (Nr.~9–11)} \verztrenn \verz{e2}{2 Nullstellen (Nr.~12–14)} \verztrenn \verz{e3}{3 Einheit 3 (Nr.~15–16)} \verztrenn \verz{e4}{4 Symmetrie (Nr.~17–18)} \verztrenn \verz{e5}{5 Transformationen (Nr.~19–20)} \verztrenn \verz{e6}{6 Graph und Term (Nr.~21–22)} \verztrenn \verz{abhaken}{Das kann ich}}
% Zone „Das kennst du schon“ – aus bank/funktionsklassen-und-eigenschaften/zone.jsonl (zusammenbau v0.1)
\einheitenkopf[zone][Kennst du schon]{Das kennst du schon}
\setcounter{aufgabe}{0}

%% TODO Titel: Ich-kann-Satz fehlt in der Bank (Platzhalter: Kettenname) – Nr. 1
%% TODO Anweisung: ein Satz über der ersten Teilaufgabe fehlt in der Bank – Nr. 1
\begin{aufgabe}{Termwerte berechnen}
\begin{teile}
\teil $f(x) = x^2 + 4$ – wie groß ist $f(3)$? f(3) = \leerfeld
\teil $h(x) = \frac{1}{2}x^2 - 3x$ – wie groß ist $h(-4)$? h(−4) = \leerfeld
\end{teile}
\end{aufgabe}

%% TODO Titel: Ich-kann-Satz fehlt in der Bank (Platzhalter: Kettenname) – Nr. 2
%% TODO Anweisung: ein Satz über der ersten Teilaufgabe fehlt in der Bank – Nr. 2
\begin{aufgabe}{Lineare Funktionen}
\begin{teile}
\teil $y = 3x - 6$ – Nullstelle? x = \leerfeld
\teil Gerade durch $A(-2 | 1)$ und $B(2 | 3)$ – Steigung? m = \leerfeld
\end{teile}
\end{aufgabe}

%% TODO Titel: Ich-kann-Satz fehlt in der Bank (Platzhalter: Kettenname) – Nr. 3
%% TODO Anweisung: ein Satz über der ersten Teilaufgabe fehlt in der Bank – Nr. 3
\begin{aufgabe}{Maßstab und Einheiten}
\begin{teile}
\teil 1 Längeneinheit entspricht 2 m. 6 Längeneinheiten – wie viele Meter? \leerfeld[m]
\teil 1 Längeneinheit entspricht 0,5 m. 3,4 Längeneinheiten – wie viele Zentimeter? \leerfeld[cm]
\end{teile}
\end{aufgabe}

%% TODO Titel: Ich-kann-Satz fehlt in der Bank (Platzhalter: Kettenname) – Nr. 4
%% TODO Anweisung: ein Satz über der ersten Teilaufgabe fehlt in der Bank – Nr. 4
\begin{aufgabe}{Quadratische Gleichungen lösen}
\begin{gleichungsraster}[2]
\gl{\text{$x^2 = 49$ – Lösungen?}} &
\gl{\text{$x^2 - 2x - 8 = 0$ – Lösungen?}} \\
\end{gleichungsraster}
\end{aufgabe}

%% TODO Titel: Ich-kann-Satz fehlt in der Bank (Platzhalter: Kettenname) – Nr. 5
%% TODO Anweisung: ein Satz über der ersten Teilaufgabe fehlt in der Bank – Nr. 5
\begin{aufgabe}{Die Grundgraphen kennen}
\begin{teile}
\teil Wie ist die Parabel $y = -x^2$ geöffnet? \\ \kreuz{nach oben} \\ \kreuz{nach unten}
\teil Welche Werte nimmt $y = \mathrm{sin}(x)$ an? \\ \kreuz{alle reellen Zahlen} \\ \kreuz{alle Zahlen von $-1$ bis $1$} \\ \kreuz{nur Zahlen größer als null}
\end{teile}
\end{aufgabe}

%% TODO Titel: Ich-kann-Satz fehlt in der Bank (Platzhalter: Kettenname) – Nr. 6
%% TODO Anweisung: ein Satz über der ersten Teilaufgabe fehlt in der Bank – Nr. 6
\begin{aufgabe}{Spiegeln im Koordinatensystem}
\begin{teile}
\teil $P(3 | 5)$ an der x-Achse gespiegelt – Bildpunkt? P'( \leerfeld | \leerfeld )
\teil $R(1{,}5 | -3)$ an der y-Achse gespiegelt – Bildpunkt? R'( \leerfeld | \leerfeld )
\end{teile}
\end{aufgabe}

%% TODO Titel: Ich-kann-Satz fehlt in der Bank (Platzhalter: Kettenname) – Nr. 7
%% TODO Anweisung: ein Satz über der ersten Teilaufgabe fehlt in der Bank – Nr. 7
\begin{aufgabe}{Termwerte berechnen – Fehler finden}
\begin{teile}
\teil Ich finde den Fehler: $r(x) = x^3 - 6x$ an der Stelle $-3$. \rechnung{r(-3) &= (-3)^3 - 6 \cdot (-3) \\ &= 27 + 18 \\ &= 45}
\end{teile}
\end{aufgabe}

%% TODO Titel: Ich-kann-Satz fehlt in der Bank (Platzhalter: Kettenname) – Nr. 8
%% TODO Anweisung: ein Satz über der ersten Teilaufgabe fehlt in der Bank – Nr. 8
\begin{aufgabe}{Termwerte berechnen – selbst rechnen}
\begin{teile}
\teil $s(x) = x^3 - 5x$ – wie groß ist $s(-2)$? s(−2) = \leerfeld
\end{teile}
\end{aufgabe}

\clearpage
% Einheit 1 von 6 – Funktionswert und Punkt (bank/funktionsklassen-und-eigenschaften/e1.jsonl, zusammenbau v0.1)
%% TODO Zeitmarke Einheit 1: Marken-Zeile nicht lesbar
%% TODO Zweigzeile Teil 1: was hier gelernt wird (Satz in Schülersprache aus dem Einheitstitel) – Einheit 1
\einheitenkopf[e1]{Einheit 1 von 6 · Funktionswert und Punkt}
\zweigzeile{Abitur GK}
\setcounter{aufgabe}{8}

%% TODO Titel: Ich-kann-Satz fehlt in der Bank (Platzhalter: Kettenname) – Nr. 9
%% TODO Anweisung: ein Satz über der ersten Teilaufgabe fehlt in der Bank – Nr. 9
\begin{aufgabe}{Funktionswert und Punkt}
%% TODO \rechenplatz{n}: Schrittzahl des Grundfalls fehlt in der Bank (nur bei fester Schreibform, 2.3 b)
\begin{teile}
\teil Wie hoch ist der Ball $2$ s nach dem Abwurf? – Wert oder Stelle? \\ \kreuz{Stelle gegeben, Wert gesucht} \\ \kreuz{Wert gegeben, Stelle gesucht}
\teil $f(x) = x^3 - 3x^2 + 2$, Stelle $2$ – Funktionswert und Punkt? f(2) = \leerfeld , P( \leerfeld | \leerfeld )
\teil $f(x) = x^3 - 2x^2 - 5$ – liegt $P(3 | 4)$ auf dem Graphen?
\teil $f(x) = x^3 + 2x^2 - 4x + 6$ – Schnittpunkt mit der y-Achse und Steigung dort? S( \leerfeld | \leerfeld ), f'(0) = \leerfeld
\teil Tee kühlt ab: $T(t) = 20 + 70 \cdot e^{-0{,}1t}$ ($T$ in °C, $t$ in Minuten nach dem Aufgießen). Anfangstemperatur und Temperatur nach $10$ Minuten (auf eine Stelle)?
\teil Lies am Graphen ab: An welchen Stellen hat $f$ den Wert $3$?

\begin{ksys}[xmin=-4,xmax=4,ymin=-2,ymax=6,ablesen]\funktion{0.5*\x^2-1.5}{f}\end{ksys}
\teil $f(x) = x^2 + 1$ und $g(x) = 2x - 3$ – senkrechter Abstand der Graphenpunkte an der Stelle $3$?
\teil Modell $T(t) = 18 + 2t$ ($T$ in °C, $t$ in Stunden ab 9 Uhr); gemessen um 12 Uhr: $25$ °C. Prozentuale Abweichung des Modellwerts vom Messwert?
\teil $g(x) = -0{,}5x + 3$ – liegt $A(4 | 1)$ auf der Geraden? Zeichne die Gerade und $A$ ein.

\begin{ksys}[xmin=-2,xmax=6,ymin=-2,ymax=5]\end{ksys}
\teil $f(x) = -x^3 + 2x^2 + 2x + 3$. Liegt $P(-2 | 13)$ auf dem Graphen? Zeige, dass $3$ eine Nullstelle ist und dass $f$ keine weitere hat \hfill (FHR 2025).
\end{teile}
\end{aufgabe}

%% TODO Titel: Ich-kann-Satz fehlt in der Bank (Platzhalter: Kettenname) – Nr. 10
%% TODO Anweisung: ein Satz über der ersten Teilaufgabe fehlt in der Bank – Nr. 10
\begin{aufgabe}{Nullstellen und Werte: Werte a mit vorgegebener Lösungsanzahl von f(x) = a über die lokalen Extremwerte am Graphen angeben}
\begin{teile}
\teil Lies am Graphen ab: Für welche Werte $a$ hat $f(x) = a$ genau drei Lösungen?

\begin{ksys}[xmin=-5,xmax=5,ymin=-20,ymax=20,ystep=4,ablesen]\funktion{\x^3-12*\x}{f}\end{ksys}
\end{teile}
\end{aufgabe}

%% TODO Titel: Ich-kann-Satz fehlt in der Bank (Platzhalter: Kettenname) – Nr. 11
%% TODO Anweisung: ein Satz über der ersten Teilaufgabe fehlt in der Bank – Nr. 11
\begin{aufgabe}{Funktionswert und Punkt – Fehler finden, Begründen, Anwendung, Darstellungswechsel}
\begin{teile}
\teil Ich finde den Fehler: Gesucht ist die Stelle, an der $f(x) = 2x + 4$ den Wert $10$ hat. \rechnung{f(10) &= 2 \cdot 10 + 4 = 24}
\teil Begründe: Liegt $P(a | b)$ auf dem Graphen von $f$, dann gilt $f(a) = b$.
\teil Ein Handwerker berechnet seine Kosten mit $K(x) = 45x + 30$ ($K$ in €, $x$ Arbeitsstunden). Ein Kunde hat $250$ € eingeplant. Reicht das für $5$ Stunden?
\teil Der Graph zeigt $f(x) = -x^2 + 4x$. Lies $f(1)$ ab und bestätige den Wert am Term.

\begin{ksys}[xmin=-1,xmax=5,ymin=-2,ymax=5,ablesen]\funktion{-\x^2+4*\x}{f}\end{ksys}
\end{teile}
\end{aufgabe}

\clearpage
% Einheit 2 von 6 – Nullstellen (bank/funktionsklassen-und-eigenschaften/e2.jsonl, zusammenbau v0.1)
%% TODO Zeitmarke Einheit 2: Marken-Zeile nicht lesbar
%% TODO Zweigzeile Teil 1: was hier gelernt wird (Satz in Schülersprache aus dem Einheitstitel) – Einheit 2
\einheitenkopf[e2]{Einheit 2 von 6 · Nullstellen}
\zweigzeile{Abitur GK}
\setcounter{aufgabe}{11}

%% TODO Titel: Ich-kann-Satz fehlt in der Bank (Platzhalter: Kettenname) – Nr. 12
%% TODO Anweisung: ein Satz über der ersten Teilaufgabe fehlt in der Bank – Nr. 12
\begin{aufgabe}{Nullstellen}
%% TODO \rechenplatz{n}: Schrittzahl des Grundfalls fehlt in der Bank (nur bei fester Schreibform, 2.3 b)
\begin{teile}
\teil $f(x) = (x + 4)(x - 7)$ – welches Verfahren? \\ \kreuz{Faktoren null setzen} \\ \kreuz{ausklammern} \\ \kreuz{substituieren} \\ \kreuz{Polynomdivision} \\ \kreuz{Lösungsformel}
\teil $f(x) = (x - 2)(x + 5)$ – Nullstellen? x1 = \leerfeld , x2 = \leerfeld
\end{teile}
\begin{gleichungsraster}[2]
\gl{\text{$f(x) = x^3 - 9x$ – Nullstellen?}} &
\gl{\text{$f(x) = x^2 - 6x + 5$ – Nullstellen?}} \\
\gl{\text{$f(x) = x^4 - 13x^2 + 36$ – Nullstellen?}} &
\gl{\text{$f(x) = x^3 - 2x^2 - 5x + 6$ – bestätige die Nullstelle $1$ und bestimme alle Nullstellen.}} \\
\end{gleichungsraster}
\begin{teile}
\teil $f(x) = (x^2 - 4) \cdot e^x$ – Nullstellen?
\teil $f(x) = x^3 - x^2 - 6x$ – alle Schnittpunkte mit den Achsen?
\teil Ein Tunnelquerschnitt wird durch $f(x) = -0{,}5x^2 + 4{,}5$ beschrieben (1 Längeneinheit entspricht 1 m, Boden auf der x-Achse). Breite des Tunnels am Boden?
\end{teile}
\begin{gleichungsraster}[2]
\gl{\text{$f(x) = 0{,}2x^4 - 2x^2 + 1{,}8$ – berechne die Koordinaten aller Schnittpunkte des Graphen mit den Koordinatenachsen (FHR 2026).}} & \\
\end{gleichungsraster}
\end{aufgabe}

%% TODO Titel: Ich-kann-Satz fehlt in der Bank (Platzhalter: Kettenname) – Nr. 13
%% TODO Anweisung: ein Satz über der ersten Teilaufgabe fehlt in der Bank – Nr. 13
\begin{aufgabe}{Nullstellen und Werte: y-Achsenschnittpunkt angeben und Anzahl der Nullstellen aus Grenzverhalten und Tiefpunkt ohne Rechnung begründen}
\begin{teile}
\teil $f(x) = x^4 - 4x + 5$ hat genau einen Tiefpunkt, $T(1 | 2)$. Schnittpunkt mit der y-Achse? Wie viele Nullstellen hat $f$? Begründe ohne Rechnung.
\end{teile}
\end{aufgabe}

%% TODO Titel: Ich-kann-Satz fehlt in der Bank (Platzhalter: Kettenname) – Nr. 14
%% TODO Anweisung: ein Satz über der ersten Teilaufgabe fehlt in der Bank – Nr. 14
\begin{aufgabe}{Nullstellen – Fehler finden, Begründen, Anwendung}
\begin{teile}
\teil Ich finde den Fehler: \rechnung{x^3 - 16x &= 0 &&\mid :x \\ x^2 &= 16 \\ x &= \pm 4}
\teil Begründe: $(x - 5)(x + 1) = 0$ gilt genau für $x = 5$ und $x = -1$.
\teil Beim Kugelstoßen gilt $h(x) = -0{,}05x^2 + 0{,}8x + 1{,}8$ ($h$ Höhe in m, $x$ waagerechte Entfernung vom Abstoß in m). Wie weit fliegt die Kugel?
\end{teile}
\end{aufgabe}

\clearpage
% Einheit 3 von 6 – Einheit 3 (bank/funktionsklassen-und-eigenschaften/e3.jsonl, zusammenbau v0.1)
%% TODO Prüfungswort Einheit 3: nicht in der Marken-Zeile
%% TODO Zeitmarke Einheit 3: Marken-Zeile nicht lesbar
%% TODO Zweigzeile Teil 1: was hier gelernt wird (Satz in Schülersprache aus dem Einheitstitel) – Einheit 3
%% TODO Einheitentitel 3 nicht in der Mappe lesbar
\einheitenkopf[e3]{Einheit 3 von 6 · Einheit 3}
\setcounter{aufgabe}{14}

%% TODO Titel: Ich-kann-Satz fehlt in der Bank (Platzhalter: Kettenname) – Nr. 15
%% TODO Anweisung: ein Satz über der ersten Teilaufgabe fehlt in der Bank – Nr. 15
\begin{aufgabe}{Definitionsbereich, Wertemenge und Schranken}
%% TODO \rechenplatz{n}: Schrittzahl des Grundfalls fehlt in der Bank (nur bei fester Schreibform, 2.3 b)
\begin{teile}
\teil $f(x) = \mathrm{ln}(x - 3)$ – größtmöglicher Definitionsbereich? D = \leerfeld
\teil $f(x) = e^x + 2$ – Wertemenge? W = \leerfeld
\teil Begründe: $f(x) = 5 - 2e^{-x}$ bleibt stets unter $5$.
\teil $f(x) = e^{-x^2}$ – Wertemenge? W = \leerfeld
\teil $f(x) = 2 + \mathrm{sin}(x)$ für $0 \le x < 2\pi$ – Wertemenge? Wie oft wird der Wert $2$ angenommen?
\teil Schar $g_k(x) = \mathrm{ln}(k \cdot x^4 + 1)$ mit $k > 0$. Gib den Definitionsbereich an, zeige, dass alle Graphen durch den Ursprung gehen, und bestimme den exakten Wert von $k$ mit $g_k(1) = 3$ \hfill (Abitur 2017 LK).
\end{teile}
\end{aufgabe}

%% TODO Titel: Ich-kann-Satz fehlt in der Bank (Platzhalter: Kettenname) – Nr. 16
%% TODO Anweisung: ein Satz über der ersten Teilaufgabe fehlt in der Bank – Nr. 16
\begin{aufgabe}{Definitionsbereich, Wertemenge und Schranken – Fehler finden, Begründen, Anwendung}
\begin{teile}
\teil Ich finde den Fehler: $f(x) = 4 - e^{-x}$ hat die Wertemenge $W = ]-\infty; 4]$.
\teil Begründe: $\mathrm{ln}(x)$ ist für $x \le 0$ nicht definiert.
\teil Ein Kühlschrank kühlt nach dem Einschalten: $T(t) = 4 + 16e^{-0{,}5t}$ ($T$ in °C, $t$ in Stunden). Lebensmittel sollen erst bei höchstens $5$ °C hinein. Ab wann geht das? Werden jemals $4$ °C erreicht?
\end{teile}
\end{aufgabe}

\clearpage
% Einheit 4 von 6 – Symmetrie (bank/funktionsklassen-und-eigenschaften/e4.jsonl, zusammenbau v0.1)
%% TODO Zeitmarke Einheit 4: Marken-Zeile nicht lesbar
%% TODO Zweigzeile Teil 1: was hier gelernt wird (Satz in Schülersprache aus dem Einheitstitel) – Einheit 4
\einheitenkopf[e4]{Einheit 4 von 6 · Symmetrie}
\zweigzeile{Abitur GK}
\setcounter{aufgabe}{16}

%% TODO Titel: Ich-kann-Satz fehlt in der Bank (Platzhalter: Kettenname) – Nr. 17
%% TODO Anweisung: ein Satz über der ersten Teilaufgabe fehlt in der Bank – Nr. 17
\begin{aufgabe}{Symmetrie}
%% TODO \rechenplatz{n}: Schrittzahl des Grundfalls fehlt in der Bank (nur bei fester Schreibform, 2.3 b)
\begin{teile}
\teil $f(x) = 2x^6 - x^2 + 5$ – gerade oder ungerade Exponenten? \\ \kreuz{achsensymmetrisch} \\ \kreuz{punktsymmetrisch} \\ \kreuz{keines von beiden}
\teil $f(x) = x^4 - 6x^2 + 2$ – Symmetrie? Begründe am Term.
\teil $f(x) = x^2 \cdot e^{-x^2}$ – Symmetrie? Begründe über $f(-x)$.
\teil $f_a(x) = a x^3 - x + a$ – Symmetrie von $f_0$?
\teil $f(x) = x^4 - 8x^2 + 7$ – Symmetrie mit Begründung und Schnittpunkt mit der y-Achse?
\teil $f$ ist achsensymmetrisch zur y-Achse, hat die Nullstelle $2{,}5$ und den Hochpunkt $H(1{,}5 | 4)$. Weitere Nullstelle und weiterer Hochpunkt?
\teil Ein symmetrischer Brückenbogen ist am Boden $24$ m breit und $6$ m hoch. Das Koordinatensystem wird so gelegt, dass die y-Achse die Symmetrieachse ist und der Boden auf der x-Achse liegt (1 Längeneinheit entspricht 1 m). Koordinaten der Fußpunkte und des höchsten Punkts?
\teil Die Gerade $y = c$ schneidet den Graphen von $f(x) = 0{,}5x^2$ in zwei Punkten, die $6$ Längeneinheiten auseinanderliegen. Wert von $c$? c = \leerfeld
\teil $f(x) = 0{,}5x^4 - 4{,}5x^2 + 2$. Begründe, warum der Graph symmetrisch zur y-Achse ist, und gib das Verhalten der Funktionswerte für $x \to -\infty$ und $x \to +\infty$ an \hfill (FHR 2026).
\end{teile}
\end{aufgabe}

%% TODO Titel: Ich-kann-Satz fehlt in der Bank (Platzhalter: Kettenname) – Nr. 18
%% TODO Anweisung: ein Satz über der ersten Teilaufgabe fehlt in der Bank – Nr. 18
\begin{aufgabe}{Symmetrie – Fehler finden, Begründen}
\begin{teile}
\teil Ich finde den Fehler: „$f(x) = x^4 - 5x^2 + 5$ ist nicht symmetrisch, weil die $5$ kein $x$ hat und deshalb ungerade zählt.“
\teil Begründe: Das Absolutglied zählt bei der Symmetrie als gerader Exponent.
\end{teile}
\end{aufgabe}

\clearpage
% Einheit 5 von 6 – Transformationen (bank/funktionsklassen-und-eigenschaften/e5.jsonl, zusammenbau v0.1)
%% TODO Zeitmarke Einheit 5: Marken-Zeile nicht lesbar
%% TODO Zweigzeile Teil 1: was hier gelernt wird (Satz in Schülersprache aus dem Einheitstitel) – Einheit 5
\einheitenkopf[e5]{Einheit 5 von 6 · Transformationen}
\zweigzeile{Abitur GK}
\setcounter{aufgabe}{18}

%% TODO Titel: Ich-kann-Satz fehlt in der Bank (Platzhalter: Kettenname) – Nr. 19
%% TODO Anweisung: ein Satz über der ersten Teilaufgabe fehlt in der Bank – Nr. 19
\begin{aufgabe}{Transformationen}
%% TODO \rechenplatz{n}: Schrittzahl des Grundfalls fehlt in der Bank (nur bei fester Schreibform, 2.3 b)
\begin{teile}
\teil $f(x)$ und $f(x - 4)$ – innen oder außen? \\ \kreuz{innen: x-Richtung} \\ \kreuz{außen: y-Richtung}
\teil $g(x) = f(x) + 3$ – wie entsteht der Graph von $g$ aus dem von $f$?
\teil $f(x) = x^2$ und $g(x) = (x + 2)^2$ – wie geht der Graph von $g$ aus dem von $f$ hervor?
\teil Der Graph von $f(x) = x^2$ ist gezeichnet. Skizziere den Graphen von $g(x) = (x - 1)^2 - 2$.

\begin{ksys}[xmin=-4,xmax=4,ymin=-3,ymax=5]\parabel{1}{0}{0}{f}\end{ksys}
\teil $f(x) = x^3 - 12x$ hat den Hochpunkt $H(-2 | 16)$. Der Graph wird um $3$ nach rechts und um $5$ nach oben verschoben. Term und Hochpunkt des Bildes?
\teil Der Graph von $g$ entsteht aus dem von $f$ durch Streckung in y-Richtung; aus $P(2 | 3)$ wird $P'(2 | 7{,}5)$. Streckfaktor?
\teil $f(x) = x^3$ und $g(x) = f(x) + 3$. Aussage: Beide Graphen haben an der Stelle $2$ dieselbe Steigung. Stimmt das?
\teil $f(x) = 3\,\mathrm{sin}(2x) + 1$ – Amplitude, Periode, Mittellinie?
\teil Die obere Randlinie eines Schildes ist $f(x) = -0{,}1x^2 + 2$ für $-4 \le x \le 4$ (1 Längeneinheit entspricht 1 dm). Die untere Randlinie entsteht durch Spiegelung an der x-Achse. Term der unteren Randlinie und Höhe des Schildes in der Mitte?
\teil $g(x) = f(x - 2) + 1$ und $f'(3) = 4$. Wie groß ist $g'(5)$?
\teil $f$ ist in $\mathbb{R}$ definiert und hat die Wertemenge $[-1; 4[$; $h(x) = -3 \cdot f(x + 2) + 2$. Gib die Wertemenge von $h$ an und begründe \hfill (Abitur 2026 LK).
\end{teile}
\end{aufgabe}

%% TODO Titel: Ich-kann-Satz fehlt in der Bank (Platzhalter: Kettenname) – Nr. 20
%% TODO Anweisung: ein Satz über der ersten Teilaufgabe fehlt in der Bank – Nr. 20
\begin{aufgabe}{Transformationen – Fehler finden, Begründen, Darstellungswechsel, Anwendung}
\begin{teile}
\teil Ich finde den Fehler: „Der Graph von $g(x) = f(x + 3)$ ist der um $3$ nach rechts verschobene Graph von $f$.“
\teil Begründe: Der Graph von $g(x) = f(x - 2)$ ist um $2$ nach rechts verschoben, obwohl im Term ein Minus steht.
\teil Die Graphen von $f(x) = x^2$ und $g$ sind gezeichnet. Lies ab: Wie entsteht $g$ aus $f$? Gib den Term von $g$ an.

\begin{ksys}[xmin=-3,xmax=5,ymin=-2,ymax=5,ablesen]\parabel{1}{0}{0}{f}\parabel{1}{2}{-1}{g}\end{ksys}
\teil Wasserstand in einem Hafen: $h(t) = 2\,\mathrm{cos}(0{,}5t) + 5$ ($h$ in m, $t$ in Stunden). Was bedeutet die Periode? Höchster und niedrigster Wasserstand?
\end{teile}
\end{aufgabe}

\clearpage
% Einheit 6 von 6 – Graph und Term (bank/funktionsklassen-und-eigenschaften/e6.jsonl, zusammenbau v0.1)
%% TODO Zeitmarke Einheit 6: Marken-Zeile nicht lesbar
%% TODO Zweigzeile Teil 1: was hier gelernt wird (Satz in Schülersprache aus dem Einheitstitel) – Einheit 6
\einheitenkopf[e6]{Einheit 6 von 6 · Graph und Term}
\zweigzeile{Abitur GK}
\setcounter{aufgabe}{20}

%% TODO Titel: Ich-kann-Satz fehlt in der Bank (Platzhalter: Kettenname) – Nr. 21
%% TODO Anweisung: ein Satz über der ersten Teilaufgabe fehlt in der Bank – Nr. 21
\begin{aufgabe}{Graph und Term}
%% TODO \rechenplatz{n}: Schrittzahl des Grundfalls fehlt in der Bank (nur bei fester Schreibform, 2.3 b)
\begin{teile}
\teil $f(x) = -2x^4 + x$ – wohin verläuft der Graph für $x \to +\infty$? \\ \kreuz{nach oben} \\ \kreuz{nach unten}
\teil Fülle die Wertetabelle von $f(x) = x^2 - 2$ aus und übertrage die Punkte ins Koordinatensystem.

\wertetabelle{x}{f(x)}{-2,-1,0,1,2}\begin{ksys}[xmin=-3,xmax=3,ymin=-3,ymax=3]\end{ksys}
\teil $f(x) = -x^3 + 3x$ hat den Tiefpunkt $(-1 | -2)$, den Hochpunkt $(1 | 2)$ und die Nullstellen $0$ und $\pm\sqrt{3}$. Berechne die Randwerte und zeichne den Graphen für $-2 \le x \le 2{,}5$.

\begin{ksys}[xmin=-3,xmax=3,ymin=-9,ymax=3]\end{ksys}
\teil Der Graph von $f(x) = x^3 - 12x$ soll für $-4 \le x \le 4$ gezeichnet werden; die Werte reichen von $-16$ bis $16$, für die Hochachse sind $8$ cm Platz. Wie viele Einheiten entspricht $1$ cm?
\teil Fischbestand in einem Teich: $B(t) = 200 \cdot 1{,}5^t$ ($t$ in Jahren). Lege eine Wertetabelle für $0 \le t \le 4$ an und zeichne den Graphen in ein selbst eingeteiltes Koordinatensystem.

\wertetabelleleer{t}{B(t)}{5}
\teil Welcher Graph gehört zu $f(x) = -x^3 + 2x^2$? Begründe über Merkmale. \\ \kreuz{A} \\ \kreuz{B} \\ \kreuz{C}

\begin{ksys}[xmin=-2,xmax=3,ymin=-3,ymax=3]\funktion{\x^3-2*\x^2}{A}\funktion{-\x^3+2*\x^2}{B}\funktion{-\x^3+2*\x}{C}\end{ksys}
\teil Eine symmetrische Schale ist oben $8$ cm breit, ihr tiefster Punkt liegt $2$ cm unter dem Rand. In einer Zeichnung ohne Achsen gehört dazu der Term $f(x) = 0{,}125x^2 - 2$ (1 Längeneinheit entspricht 1 cm). Wo liegen die Achsen?
\teil Lies am Graphen $f(3)$ ab und entscheide mit Begründung, ob $f'(3)$ positiv oder negativ ist.

\begin{ksys}[xmin=-1,xmax=5,ymin=-2,ymax=5,ablesen]\funktion{-0.5*\x^2+2*\x+1}{f}\end{ksys}
\teil Der Graph einer ganzrationalen Funktion ist gezeichnet. Welchen Grad hat sie mindestens?

\begin{ksys}[xmin=-4,xmax=4,ymin=-4,ymax=3]\funktion{0.25*\x^4-2*\x^2+1}{f}\end{ksys}
\teil $f(x) = 2\,\mathrm{cos}(x)$ – auf welcher Geraden liegen alle Hochpunkte?
\teil $f(x) = x^4 - 18x^2$ hat die Tiefpunkte $T_1(-3 | -81)$ und $T_2(3 | -81)$ und den Hochpunkt $H(0 | 0)$. Zeichne den Graphen für $-4 \le x \le 4$ mit geeigneter Achseneinteilung, insbesondere für die y-Achse \hfill (FHR 2023).

\wertetabelleleer{x}{f(x)}{9}
\end{teile}
\end{aufgabe}

%% TODO Titel: Ich-kann-Satz fehlt in der Bank (Platzhalter: Kettenname) – Nr. 22
%% TODO Anweisung: ein Satz über der ersten Teilaufgabe fehlt in der Bank – Nr. 22
\begin{aufgabe}{Graph und Term – Fehler finden, Begründen, Darstellungswechsel, Anwendung}
\begin{teile}
\teil Ich finde den Fehler: „Der Graph hat zwei Nullstellen und drei Extrempunkte, also ist die Funktion vom Grad $2$.“
\teil Begründe: Eine ganzrationale Funktion mit drei Extrempunkten hat mindestens den Grad vier.
\teil Lies am Graphen die Werte an den Stellen $-2$, $0$ und $2$ ab und trage sie in die Tabelle ein.

\begin{ksys}[xmin=-3,xmax=3,ymin=-2,ymax=3,ablesen]\funktion{0.5*\x^2-1}{f}\end{ksys}\wertetabelle{x}{f(x)}{-2,0,2}
\teil Der Graph zeigt die Zahl $B$ der Besucher eines Freizeitparks in Tausend, $t$ in Stunden nach 8 Uhr. Lies den Hochpunkt ab und deute beide Koordinaten.

\begin{ksys}[xmin=0,xmax=12,ymin=0,ymax=10,xstep=1,ystep=1,xlabel=t in h,ylabel=B in Tausend,ablesen]\funktion{-0.25*\x^2+3*\x}{B}\end{ksys}
\end{teile}
\end{aufgabe}

% Abhakseite „Das kann ich“ – Zone nur im Gesamt (\mitzone)
%% TODO Ich-kann-Sätze: Platzhalter wie in den Aufgabendateien
\begin{abhakseite}
\ifdefined\mitzone
\abhakgruppe{Kennst du schon}
\abhak{1}{Termwerte berechnen}
\abhak{2}{Lineare Funktionen}
\abhak{3}{Maßstab und Einheiten}
\abhak{4}{Quadratische Gleichungen lösen}
\abhak{5}{Die Grundgraphen kennen}
\abhak{6}{Spiegeln im Koordinatensystem}
\abhak{7}{Termwerte berechnen – Fehler finden}
\abhak{8}{Termwerte berechnen – selbst rechnen}
\fi
\abhakgruppe{Einheit 1 · Funktionswert und Punkt}
\abhak{9}{Funktionswert und Punkt}
\abhak{10}{Nullstellen und Werte: Werte a mit vorgegebener Lösungsanzahl von f(x) = a über die lokalen Extremwerte am Graphen angeben}
\abhak{11}{Funktionswert und Punkt – Fehler finden, Begründen, Anwendung, Darstellungswechsel}
\abhakgruppe{Einheit 2 · Nullstellen}
\abhak{12}{Nullstellen}
\abhak{13}{Nullstellen und Werte: y-Achsenschnittpunkt angeben und Anzahl der Nullstellen aus Grenzverhalten und Tiefpunkt ohne Rechnung begründen}
\abhak{14}{Nullstellen – Fehler finden, Begründen, Anwendung}
\abhakgruppe{Einheit 3 · Einheit 3}
\abhak{15}{Definitionsbereich, Wertemenge und Schranken}
\abhak{16}{Definitionsbereich, Wertemenge und Schranken – Fehler finden, Begründen, Anwendung}
\abhakgruppe{Einheit 4 · Symmetrie}
\abhak{17}{Symmetrie}
\abhak{18}{Symmetrie – Fehler finden, Begründen}
\abhakgruppe{Einheit 5 · Transformationen}
\abhak{19}{Transformationen}
\abhak{20}{Transformationen – Fehler finden, Begründen, Darstellungswechsel, Anwendung}
\abhakgruppe{Einheit 6 · Graph und Term}
\abhak{21}{Graph und Term}
\abhak{22}{Graph und Term – Fehler finden, Begründen, Darstellungswechsel, Anwendung}
\end{abhakseite}

\end{document}
